% Proposed replacement for the supplied theory subsection.
% Requires amsmath, amssymb, and natbib. Retains the supplied citation keys.
% Target length: half to one manuscript page, excluding the figure.
% See theory_audit_2026-10-07.txt before selecting the final figure.
\subsection{Target-dependent fitting in a tractable model}
\label{sec:linear_theory}

Can changing what a model predicts change how quickly it fits a fixed dataset,
even when the learning dynamics can be solved exactly? We begin with linear
prediction and then add fixed nonlinear features, while training only the
final linear readout. Both models remain exactly solvable at one fixed flow
time. They isolate sources of target-dependent fitting without modeling the
full generation process.

Draw $\mathcal D=\{\mathbf x_i\}_{i=1}^n$ independently from
$\mathcal N(0,I_d)$. Reuse these clean examples, but draw fresh
$\boldsymbol\epsilon\sim\mathcal N(0,I_d)$ in
$\mathbf z=t\mathbf x+(1-t)\boldsymbol\epsilon$.
Here $t\in(0,1)$ is fixed and $\tau$ denotes optimization time.
For both models, train $f_q(\mathbf z;\tau)=B_q(\tau)\boldsymbol\phi(\mathbf z)$
by gradient flow on the per-coordinate squared error, with targets
$\mathbf y_x=\mathbf x$, $\mathbf y_\epsilon=\boldsymbol\epsilon$, and
$\mathbf y_v=\mathbf v=\mathbf x-\boldsymbol\epsilon$.
Here $q\in\{x,\epsilon,v\}$ and the feature map is
\[
 \boldsymbol\phi(\mathbf z)=
 \begin{cases}
 \mathbf z,&\text{linear model},\\
 \tanh(W\mathbf z/\sqrt d),&\text{random-feature network}.
 \end{cases}
\]
In the second case, $W\in\mathbb R^{p\times d}$ has independent standard
Gaussian entries and is frozen. Only $B_q$ is trained.
All readouts start at zero and use the same learning-rate normalization.
For either choice, define
$U=\mathbb E_{\mathrm{tr}}[\boldsymbol\phi\boldsymbol\phi^\top]$ and
$V_q=\mathbb E_{\mathrm{tr}}[\boldsymbol\phi\mathbf y_q^\top]$,
where the training expectation averages over stored examples and fresh noise.
The loss is quadratic in $B_q$. Absorbing its common factor $2/d$ into
$\tau$ gives the exact dynamics and solution
\begin{equation}
 \dot B_q=V_q^\top-B_qU,\qquad
 B_q(\tau)=V_q^\top g_\tau(U),
 \label{eq:main_feature_solution}
\end{equation}
where $g_\tau(\lambda)=(1-e^{-\lambda\tau})/\lambda$ and $g_\tau(0)=\tau$.

\paragraph{Linear model.}
For $\boldsymbol\phi(\mathbf z)=\mathbf z$, express each prediction as a velocity
$\widehat{\mathbf v}_q=A_q\mathbf z$:
\begin{equation}
 A_x=\frac{B_x-I_d}{1-t},\qquad
 A_\epsilon=\frac{I_d-B_\epsilon}{t},\qquad A_v=B_v.
 \label{eq:main_velocity_conversion}
\end{equation}
Let $S=n^{-1}\sum_i\mathbf x_i\mathbf x_i^\top$ and
$U=t^2S+(1-t)^2I_d$.
All three velocity predictors obey
\begin{equation}
 \begin{gathered}
 \dot A_q=tS-(1-t)I_d-A_qU,
 \\
 (A_x(0),A_\epsilon(0),A_v(0))
 =\left(-\frac{I_d}{1-t},\frac{I_d}{t},0\right).
 \end{gathered}
 \label{eq:main_shared_flow}
\end{equation}
Thus the equation is shared, but zero native outputs give different starting
velocities. The trajectories converge to the same empirical solution.
Matching their initial velocities makes their entire trajectories identical.

We measure velocity mean squared error per coordinate, $R^v_{q,r}$,
either on stored clean examples ($r=\mathrm{tr}$) or fresh population
examples ($r=P$), with fresh corruption in both cases. Define
\begin{equation}
 \mathcal G_q=R^v_{q,P}-R^v_{q,\mathrm{tr}},\qquad
 \Delta\mathcal G_q(\tau)=\mathcal G_q(\tau)-\mathcal G_q(0).
 \label{eq:main_gap}
\end{equation}
A positive change means that the training--population error difference has
widened in favor of the training set. It does not require population error
to increase, and it is not a generated-copying fraction.

At $t=1/2$, the expected early changes satisfy
\begin{equation}
\begin{aligned}
 \mathbb E_{\mathcal D}\Delta\mathcal G_x(\tau)
 &=\frac{2(d+1)}{n}\tau+O(\tau^2),\\
 \mathbb E_{\mathcal D}\Delta\mathcal G_v(\tau)
 &=\frac{d+1}{2n}\tau+O(\tau^2),\\
 \mathbb E_{\mathcal D}\Delta\mathcal G_\epsilon(\tau)
 &=\frac{d+1}{16n}\tau^3+O(\tau^4).
\end{aligned}
\label{eq:main_early_gap}
\end{equation}
Under this initialization, clean prediction develops the largest initial
expected gap, while noise prediction has a delayed leading term.
The model fits fluctuations in sample second moments. At this Gaussian
midpoint, zero velocity is already the population-optimal predictor;
population error must therefore be read alongside the gap.

\paragraph{One step beyond linear prediction.}
With the frozen nonlinear feature map, the same exact readout solution
holds, as in related fixed-feature analyses
\citep{bonnaire2025dynamical,george2025denoisingscorematchingrandom}.
The feature second-moment matrix $U$ fixes shared response times; the feature--target
moments $V_q$ determine how each target is fitted.
Nonlinear features can make the noise--feature relationship depend on the
clean example already in the first update. Consequently, this tractable
extension need not preserve the linear model's ordering or cubic noise term.
Appendix~\ref{app:target_fitting} gives the full derivations and shows how
variation in feature--target relationships determines the initial expected
gap. Thus the analysis identifies mechanisms under specified initialization
and feature choices, rather than a universal ranking of prediction targets.
